%% kapa: exam=1 round=33 year=2022 period=1 form=A subject=real_estate subject_ko=부동산학원론 seq=27 number=107
%% source: Q-Net 감정평가사 기출문제 (cst00309) artlSeq=5209997 file=raw/1008/5209997/2022년도 제33회 감정평가사 1차시험 1교시(민법,경제학,부동산학).pdf
%% conversion: kordoc 4.24.51 -> tools/gen1.py (auto); review=proofread-sonnet; flags=none
\begin{question}{107}
부동산개발의 타당성분석 유형을 설명한 것이다. ( )에 들어갈 내용으로 옳게 연결된 것은?
\begin{framedtext}
○ ( ㄱ )은 부동산이 현재나 미래의 시장상황에서 매매 또는 임대될 수 있는 가능성을 분석하는 것이다. \par
○ ( ㄴ )은 개발업자가 대상부동산에 대해 수립한 사업안들 중에서 최유효이용을 달성할 수 있는 방식을 판단할 수 있도록 자료를 제공해주는 것이다. \par
○ ( ㄷ )은 주요 변수들의 초기 투입값을 변화시켜 적용함으로써 낙관적 또는 비관적인 상황에서 발생할 수 있는 수익성 및 부채상환능력 등을 예측하는 것이다.
\end{framedtext}
\begin{choices}
  \item ㄱ: 시장성분석, ㄴ: 민감도분석, ㄷ: 투자분석
  \item ㄱ: 민감도분석, ㄴ: 투자분석, ㄷ: 시장성분석
  \item ㄱ: 투자분석, ㄴ: 시장성분석, ㄷ: 민감도분석
  \item ㄱ: 시장성분석, ㄴ: 투자분석, ㄷ: 민감도분석
  \item ㄱ: 민감도분석, ㄴ: 시장성분석, ㄷ: 투자분석
\end{choices}
\end{question}
\answer{4}
